% Privacy and sufficiency notes for the FedGM project (2026-11-16).
% Self-contained; paste into Overleaf. Experiments: docs/findings.md N17 (E1), N18 (E2).
% Not compiled locally (no LaTeX on the machine): check the build in Overleaf.
\documentclass[11pt]{article}
\usepackage[margin=1in]{geometry}
\usepackage{amsmath,amssymb,amsthm}
\usepackage{hyperref}
\newtheorem{proposition}{Proposition}
\newtheorem{lemma}[proposition]{Lemma}
\newtheorem{corollary}[proposition]{Corollary}
\theoremstyle{remark}
\newtheorem{remark}[proposition]{Remark}
\newcommand{\R}{\mathbb{R}}
\newcommand{\E}{\mathbb{E}}
\DeclareMathOperator{\Var}{Var}
\DeclareMathOperator{\rowsp}{row}
\DeclareMathOperator{\range}{range}

\title{What federated matching must share, and what it reveals}
\author{FedGM project notes}
\date{16 November 2026}

\begin{document}
\maketitle

\paragraph{Scope.} Sites $k=1,\dots,K$ hold records $x\in\R^p$ (plus treatment and outcome);
$n_k$ records at site $k$, $N=\sum_k n_k$. A \emph{release} is anything a site sends. We state
(i) which sums suffice for each method's computations, (ii) what exact releases reveal
unavoidably, (iii) membership inference on released means, (iv) reconstruction from the
messages of the Fed-GenMatch draft (B.~Nestor), (v) per-patient sensitivities for a
differentially private version, and (vi) known lower bounds. Statements marked
\emph{(checked)} were verified numerically in the repository.

\section{Sufficiency: what each computation needs}

\begin{proposition}[Pooled moments from one round of sums]\label{prop:pooled}
Let site $k$ send $n_k$, $T_k=\sum_{i\in k}x_i$ and $Q_k=\sum_{i\in k}x_ix_i^\top$. Then the
pooled mean and covariance are
\[
\bar x=\frac{1}{N}\sum_k T_k,\qquad
\hat S=\frac{1}{N-1}\Big(\sum_k Q_k-N\bar x\bar x^\top\Big).
\]
If only the diagonal is needed (per-covariate scaling, as in R's \texttt{GenMatch}, which
z-scores each covariate and never uses a full covariance), the site sends $n_k$, $T_k$ and
$q_k=\sum_{i\in k}x_i\circ x_i$: $2p+1$ numbers instead of $1+p+p(p+1)/2$.
\end{proposition}
\begin{proof}
$\sum_i(x_i-\bar x)(x_i-\bar x)^\top=\sum_i x_ix_i^\top-N\bar x\bar x^\top$, and both terms
are sums over sites. The diagonal uses only $\sum_i x_{il}^2$.
\end{proof}

\begin{remark}
The average of the site covariances is not $\hat S$: it omits the between-site term
$\sum_k n_k(\bar x_k-\bar x)(\bar x_k-\bar x)^\top/(N-1)$.
\end{remark}

\begin{proposition}[Exact outputs reveal the other site's sums when $K=2$]\label{prop:twosite}
If the protocol's output includes $N$, $\bar x$ and $\hat S$ exactly, then with $K=2$ site 1
learns $(n_2,T_2,Q_2)$ exactly, whatever the protocol (secure aggregation included):
$n_2=N-n_1$, $T_2=N\bar x-T_1$, $Q_2=(N-1)\hat S+N\bar x\bar x^\top-Q_1$.
\end{proposition}
\noindent This is a lower bound on disclosure for exact computation: protecting the other site's
sums with two sites requires noise on the output itself.

\begin{proposition}[Fed-GenMatch's objective needs only matched-set sums]
For a candidate $W$, let $\mathcal M_k(W)$ be the controls of site $k$ matched by the
assignment. With the treated mean $\bar x^t$ and variances $s_t^2$ public (the draft's
single-arm setting), the draft's imbalance $\max_l|\mathrm{SMD}_l(W)|$ is a function of
$\big(n_k^*,\ \sum_{j\in\mathcal M_k}x_j,\ \sum_{j\in\mathcal M_k}x_j\circ x_j\big)_{k}$
(the draft's eq.~10). It is not a function of these sums for KS- or $p$-value-based
objectives, which need the matched controls' order statistics.
\end{proposition}

\begin{proposition}[FedDSM's fits are exact from sums]
The logistic propensity MLE by Newton--Raphson needs, at each iterate $\beta$, the sums
$\sum_i(w_i-p_i(\beta))x_i$ and $\sum_i p_i(\beta)(1-p_i(\beta))x_ix_i^\top$ (per site, then
added). The least-squares prognostic score on a basis $b(x)$, cross-fitted over folds, needs
per site and fold $\sum_i b(x_i)b(x_i)^\top$ and $\sum_i b(x_i)y_i$. Both give the pooled
estimates exactly \emph{(checked)}.
\end{proposition}

\section{Membership inference on released means}\label{sec:mia}

Let $\phi:\R^p\to\R^D$ be the released features (e.g.\ $x$ and $x\circ x$; $D=2p$), with
population mean $\mu$ and covariance $\Sigma_\phi$ (full rank on its support). A group $G$ of $n$
i.i.d.\ members is released as $\hat m=\frac1n\sum_{i\in G}\phi(x_i)$. The attacker knows the
population and a target $x$, and scores $s(x)=(\phi(x)-\mu)^\top\Sigma_\phi^{-1}(\hat m-\mu)$
(the whitened inner-product, or tracing, attack).

\begin{proposition}\label{prop:mia}
Write $u=\Sigma_\phi^{-1/2}(\phi(x)-\mu)$, so $\E u=0$, $\mathrm{Cov}\,u=I_D$, and let
$V=\Var\|u\|^2$. For a non-member, $\E s=0$ and $\Var s=D/n$. For a member,
$\E s=D/n$ and $\Var s=\{(n-1)D+V\}/n^2$. If $u\sim N(0,I_D)$ (so $V=2D$) and $s$ is
approximately normal (large $D$),
\[
\mathrm{AUROC}\approx\Phi\!\Big(\sqrt{D/n}\,\big/\sqrt{2+1/n}\Big)\approx\Phi\big(\sqrt{D/(2n)}\big),
\qquad
\mathrm{TPR}(\alpha)\approx1-\Phi\big(z_{1-\alpha}-\sqrt{D/n}\big).
\]
\end{proposition}
\begin{proof}
Non-member: $s=u^\top\bar u$ with $\bar u=\frac1n\sum_{i\in G}u_i$ independent of $u$, so
$\Var s=\E[u^\top\E(\bar u\bar u^\top)u]=\E\|u\|^2/n=D/n$. Member $x=x_1$:
$s=\frac1n\|u_1\|^2+\frac1n u_1^\top\sum_{i\ge2}u_i$; the two terms are uncorrelated, with
variances $V/n^2$ and $(n-1)D/n^2$. With normal scores, $\mathrm{AUROC}=\Phi\big(\Delta/\sqrt{\sigma_0^2+\sigma_1^2}\big)$.
\end{proof}

\begin{remark}[What the experiment shows]
$n$ and $p$ enter only through $D/n$, the released numbers per patient: $D=2p$ for means and
variances, $p+p(p+1)/2$ for the covariance, $\binom{p+4}{4}-1$ distinct moments for FedDSM's
quadratic-basis fold sums. The formula reproduces the Nov.~4 pilot (e.g.\ $D=65$, $n=35$:
0.83 and 16.8\% at 1\% FPR vs 0.85 and 16.9\%). In E1 it holds wherever $D/n\le1$; above, the
attack falls below it (up to $-0.13$ AUROC), most for non-Gaussian features (products,
fourth powers, log-normal data). Keeping the general $V$ over-corrects: the member scores are
then skewed and the normal approximation of a rank statistic fails. The cause is open.
\end{remark}

\section{Reconstruction from Fed-GenMatch's cross terms}

Notation of the draft: $S=\Sigma^{-1/2}$ symmetric (pseudo-inverse square root if $\Sigma$ is
singular), $\tilde X_t\in\R^{n_t\times p}$ the public whitened treated rows, $z_j=Sx_j$ for
control $j$, and at evaluation $g$ the cross terms $b_j^{(g)}=\tilde X_tW_gz_j\in\R^{n_t}$ with
$W_g$ diagonal, known to the coordinator. After $G$ evaluations the coordinator knows
$A_Gz_j$, $A_G=[\tilde X_tW_1;\dots;\tilde X_tW_G]$, so it knows the projection of $z_j$ on
$\rowsp(A_G)=\sum_gW_g\,U$, where $U=\rowsp(\tilde X_t)$.

\begin{lemma}[What can ever be exposed]\label{lem:exposed}
Let $L=\{l:\ u_l\neq0\text{ for some }u\in U\}$. For diagonal $W$ ranging over any set with
non-empty interior (e.g.\ the draft's box $[10^{-2},10^2]^p$),
$\operatorname{span}\{Wu:\ W,\ u\in U\}=\operatorname{span}\{e_l: l\in L\}$, and finitely many
generic $W_g$ reach it. Coordinates $l\notin L$ (those with $\tilde x_{t,il}=0$ for every treated
$i$) are never exposed, whatever the search does.
\end{lemma}
\begin{proof}
$Wu=\operatorname{diag}(u)\,w$ is linear in $w$; over a set with non-empty interior its span is
$\range\operatorname{diag}(u)=\operatorname{span}\{e_l:u_l\neq0\}$. Sum over a basis of $U$.
\end{proof}

\begin{corollary}[The draft's normaliser exposes everything at once]\label{cor:oneeval}
With $\Sigma=S_t$, $U=\range(S_t)$ and $z_j\in\range(S_t)$. At the first evaluation
($W=I$) the coordinator recovers $z_j$ exactly, hence $x_j$ modulo $\ker S_t$ (the directions
with no treated variation). No later evaluation adds anything.
\end{corollary}
\begin{proof}
$\rowsp(X_tS)=S\operatorname{span}\{x_{t,i}\}\supseteq S\range(S_t)=\range(S_t)$, and
$\rowsp(X_tS)\subseteq\range(S)=\range(S_t)$. If $\tilde X_tz=0$ with $z\in\range(S_t)$ then
$z\in\range(S_t)\cap\range(S_t)^\perp=\{0\}$.
\end{proof}
\noindent E2 (IST single-arm, 31 covariates): full exposure after one evaluation in 10/10
datasets at every cohort size; re-identification 80\% of controls at $n_t=20$, $>90\%$ from
$n_t=31$. With a full-rank (shrunk) normaliser the exposed part saturates within 5
evaluations but misses the covariates equal to 0 for all treated, as the lemma predicts.

\section{The matched-set sums}

At evaluation $g$ site $k$ sends $s_g=\sum_{j\in\mathcal M_g}x_j$ (and $n^*$, $\sum x\circ x$);
the coordinator knows $\mathcal M_g$ because it solved the assignment. With $E$ the
evaluations$\times$controls indicator matrix of site $k$, it knows $EX_k$.

\begin{proposition}[Identification]\label{prop:matchedsets}
$x_j$ is determined by the matched-set sums iff $e_j\in\rowsp(E)$; in particular if two
evaluations' matched sets differ exactly by $\{j\}$. \emph{(E2: 97--99\% of the controls ever
matched are identified, all 31 covariates; 91--92\% of all controls with the whole cohort.)}
\end{proposition}

\begin{lemma}[One control moves the matched set by at most one swap]\label{lem:swap}
Consider the optimal $1{:}m$ assignment without replacement (draft eq.~8), assumed unique,
before and after changing control $j$'s covariates (everything else fixed). Then the matched
sets $\mathcal M,\mathcal M'$ satisfy $|(\mathcal M\triangle\mathcal M')\setminus\{j\}|\le2$:
apart from $j$, at most one control leaves and one enters. Who is matched to whom can change for
many treated. \emph{(Checked on 3,000 random instances: the set difference never exceeded 2,
while up to 69\% of treated changed match.)}
\end{lemma}
\begin{proof}
Duplicate each treated row $m$ times, so the assignment is a minimum-cost matching saturating
the treated side. The symmetric difference of the old and new optimal matchings is a union of
vertex-disjoint alternating paths and cycles; every treated vertex has degree 0 or 2 in it, so
paths end at controls. A component $C$ not containing $j$ has the same edge costs in both
instances, so swapping $C$ cannot strictly improve either matching; by uniqueness $C$ is
empty. The one remaining component contains $j$; a cycle leaves the matched set unchanged, a
path changes it by its two control endpoints, one leaving and one entering.
\end{proof}

\begin{corollary}[Per-evaluation sensitivity of the matched-set sums]
Let covariate $l$ take values in an interval of length $r_l$ (and $x_l^2$ in one of length
$r^{(2)}_l$). Replacing one control changes $\sum_{\mathcal M}x_l$ by at most $2r_l$ and
$\sum_{\mathcal M}x_l^2$ by at most $2r^{(2)}_l$ (the changed control plus one swap), and the
per-site counts by at most 1 at each of at most two sites. So the moments of eq.~9 have a bounded
sensitivity, of the order of one record, under unique optima. (Ties need a fixed tie-break.)
\end{corollary}

\begin{remark}[Reconstruction barrier for subset sums]
The matched-set sums are subset-sum queries on each site's records. For binary attributes,
answers to $O(n\log^2n)$ random subset sums with error $o(\sqrt n)$ allow reconstruction of
all but a vanishing fraction of the records (Dinur \& Nissim 2003; Dwork, McSherry \& Talwar
2007). The search releases thousands of such sums per site (E2: 1,100--10,100 evaluations,
median 2,250; sites of 92--1,161 controls), on sets that differ by single swaps, which makes the system
easier to solve than random queries. Any protection of these sums therefore needs noise of
order at least $\sqrt{n_k}$ per sum, or far fewer released sums.
\end{remark}

\section{Per-patient sensitivity of the cross terms and composition}

Adjacency: replace one control, with records in a domain of $\ell_2$-diameter $\rho$ and norm at
most $R$.

\begin{proposition}\label{prop:sens}
At one evaluation with weights $W$, replacing control $j$ changes only column $b_j$ and $C_j$:
\[
\|\Delta b_j\|_2\le\|\tilde X_tWS\|_2\,\rho,\qquad
\|\Delta b_j\|_1\le\sum_{i=1}^{n_t}\|W\tilde x_{t,i}\|_2\,\|S\|_2\,\rho=\sum_i\Delta_i,\qquad
|\Delta C_j|\le\|W\|_2\|S\|_2^2R^2,
\]
where $\Delta_i$ is the draft's per-entry sensitivity (eq.~14). Hence Laplace noise of scale
$\Delta_i/\varepsilon$ on each entry, as in the draft, gives $n_t\varepsilon$-DP per evaluation
for the cross terms alone, not $\varepsilon$; adding $C_j$ and the matched-set sums adds
their own terms; over $G$ (adaptively chosen) evaluations, basic composition gives
$G(n_t\varepsilon+\varepsilon_C+\varepsilon_{\mathrm{mom}})$.
\end{proposition}
\begin{proof}
$\Delta b_j=\tilde X_tWS(x_j'-x_j)$; entry $i$ is $(SW\tilde x_{t,i})^\top(x'_j-x_j)$, bounded by
$\|SW\tilde x_{t,i}\|_2\rho\le\|W\tilde x_{t,i}\|_2\|S\|_2\rho$; sum over $i$. For $C_j=z^\top Wz$
with $z=Sx$, both values lie in $[0,\|W\|_2\|S\|_2^2R^2]$.
\end{proof}

\begin{remark}[Gaussian mechanism over the search]
Releasing each evaluation's $(b_j,C_j)$ with Gaussian noise of variance
$\sigma_g^2=\Delta_{2,g}^2/(2\varrho_g)$ is $\varrho_g$-zCDP; adaptive composition adds the
$\varrho_g$, and $\varrho$-zCDP implies $(\varrho+2\sqrt{\varrho\ln(1/\delta)},\delta)$-DP (Bun \&
Steinke 2016). For a fixed total budget split over $G$ evaluations, the noise per release grows
like $\sqrt G$: about $70\times$ the single-release noise at $G\approx5{,}000$ (the median search
length with all 235 treated, N15). E2 found that even single-release noise at the draft's calibration leaves no match
unchanged, for $\varepsilon$ up to $1{,}000$ per entry.
\end{remark}

\section{Lower bounds to build on}

\begin{enumerate}
\item \emph{Exact outputs:} Proposition~\ref{prop:twosite}: with two sites, exact pooled mean
  and covariance reveal the other site's sums; no protocol avoids it.
\item \emph{Reconstruction:} the subset-sum barrier above (Dinur \& Nissim 2003).
\item \emph{Price of privacy for shared means:} for $(\varepsilon,\delta)$-DP estimation of a
  $d$-dimensional Gaussian mean, the minimax squared error is of order
  $d/n+d^2\log(1/\delta)/(n^2\varepsilon^2)$ up to logarithmic factors (Cai, Wang \& Zhang 2021,
  \emph{Ann.\ Statist.}); covariance estimation has analogous $d^2$-type terms (Kamath et
  al.\ 2019). These bound the utility cost of privatising the statistics of
  Section~1. \emph{To check against the papers before citing.}
\end{enumerate}

\end{document}
